% GENERATED FILE. Edit canonical lessons, then run npm run generate:books.
% canonical-source-hash: fnv1a64:e7e83de2acf84a9b
% canonical-lessons: FA-C378-dastmozd, FA-C378-daramad, FA-C378-ezafe-kari, FA-C378-padash, FA-C378-mazaya

\chapter{Earning a living}
\label{ch:fa-earning-a-living}
\hlchaptermodality{378}

\begin{chapteropening}
\textbf{By the end of this chapter:} \emph{I can talk about wages and income, overtime, bonuses and the benefits of a job.}

Complete the last lesson of chapter 378: I can talk about wages and income, overtime, bonuses and the benefits of a job.
\end{chapteropening}

\section[\texorpdfstring{\fa{دستمزد} (dastmozd)}{دستمزد (dastmozd)}]{\fa{دستمزد} (dastmozd) — a wage, pay}

\label{lesson:FA-C378-dastmozd}



\begin{quote}
Before the new one: say the Persian for prevention, then the Persian for to matter.
\end{quote}

\subsection*{You'll want to know: \fa{دستمزد}}
\textbf{\fa{دستمزد}} (\emph{dastmozd}) — \textquotedblleft{}a wage, pay\textquotedblright{}.

\begin{quote}
\fa{دستمزد} \fa{کارگرها} \fa{کم} \fa{است}.
\end{quote}

\begin{quote}
\emph{dastmozd-e kârgarhâ kam ast.}
\end{quote}

\begin{quote}
The workers' wages are low.
\end{quote}

\emph{Dast} \textquotedblleft{}hand\textquotedblright{} + \emph{mozd} \textquotedblleft{}wage\textquotedblright{}. For daily or hourly pay; a monthly salary is \emph{hoquq}.

\begin{cousinweb}
\textbf{\fa{دستمزد}} is a native Persian compound built from \emph{dast} \textquotedblleft{}hand\textquotedblright{}, which EtymDB traces to PIE \emph{ǵ\textsuperscript{h}es-} \textquotedblleft{}hand\textquotedblright{}, and \emph{mozd} \textquotedblleft{}wage\textquotedblright{}, usually traced to Middle Persian \emph{mizd}: pay for the work of your hands.
\end{cousinweb}

\begin{grammarlens}[title={What you've built}]
One more word for telling, explaining and describing at B1.
\end{grammarlens}

\subsection*{Your turn}
\begin{itemize}
\raggedright
  \item \emph{Say it:} \emph{dastmozd}
  \item \emph{Say it:} \emph{dastmozd}, once more
  \item \emph{Say it:} a sentence of your own with \emph{dastmozd}
  \item \emph{Recall:} say the Persian for prevention, then the Persian for to matter, then say \emph{dastmozd} again
\end{itemize}

\begin{tcolorbox}[breakable,colback=teal!4,colframe=teal!35!black,title={Before you move on}]
What does \fa{دستمزد} mean? (\textquotedblleft{}A wage, pay\textquotedblright{}.) Say it once more.
\end{tcolorbox}
\section[\texorpdfstring{\fa{درآمد} (darâmad)}{درآمد (darâmad)}]{\fa{درآمد} (darâmad) — an income}

\label{lesson:FA-C378-daramad}



\begin{quote}
Before the new one: say the Persian for to matter, then the Persian for a wage.
\end{quote}

\subsection*{You'll want to know: \fa{درآمد}}
\textbf{\fa{درآمد}} (\emph{darâmad}) — \textquotedblleft{}an income\textquotedblright{}.

\begin{quote}
\fa{درآمد} \fa{او} \fa{از} \fa{کار} \fa{تازه} \fa{بیشتر} \fa{است}.
\end{quote}

\begin{quote}
\emph{darâmad-e u az kâr-e tâze bishtar ast.}
\end{quote}

\begin{quote}
His income from the new job is higher.
\end{quote}

\emph{Dar} \textquotedblleft{}in\textquotedblright{} + \emph{âmad} \textquotedblleft{}came\textquotedblright{}: what comes in. Plural \emph{darâmadhâ}.

\begin{cousinweb}
\textbf{\fa{درآمد}} is native Persian, built from \emph{dar} \textquotedblleft{}in\textquotedblright{} + \emph{âmad}, the past stem of \emph{âmadan} \textquotedblleft{}to come\textquotedblright{}, which EtymDB traces to PIE \emph{g\textsuperscript{w}em-} \textquotedblleft{}to step\textquotedblright{}. Income is what comes in.
\end{cousinweb}

\begin{grammarlens}[title={What you've built}]
One more word for telling, explaining and describing at B1.
\end{grammarlens}

\subsection*{Your turn}
\begin{itemize}
\raggedright
  \item \emph{Say it:} \emph{darâmad}
  \item \emph{Say it:} \emph{darâmad}, once more
  \item \emph{Say it:} a sentence of your own with \emph{darâmad}
  \item \emph{Recall:} say the Persian for to matter, then the Persian for a wage, then say \emph{darâmad} again
\end{itemize}

\begin{tcolorbox}[breakable,colback=teal!4,colframe=teal!35!black,title={Before you move on}]
What does \fa{درآمد} mean? (\textquotedblleft{}An income\textquotedblright{}.) Say it once more.
\end{tcolorbox}
\section[\texorpdfstring{\fa{اضافه} \fa{کاری} (ezâfe kâri)}{اضافه کاری (ezâfe kâri)}]{\fa{اضافه} \fa{کاری} (ezâfe kâri) — overtime}

\label{lesson:FA-C378-ezafe-kari}



\begin{quote}
Before the new one: say the Persian for a wage, then the Persian for an income.
\end{quote}

\subsection*{You'll want to know: \fa{اضافه} \fa{کاری}}
\textbf{\fa{اضافه} \fa{کاری}} (\emph{ezâfe kâri}) — \textquotedblleft{}overtime\textquotedblright{}.

\begin{quote}
\fa{این} \fa{ماه} \fa{خیلی} \fa{اضافه} \fa{کاری} \fa{کردم}.
\end{quote}

\begin{quote}
\emph{in mâh kheyli ezâfe kâri kardam.}
\end{quote}

\begin{quote}
This month I did a lot of overtime.
\end{quote}

\emph{Ezâfe} \textquotedblleft{}extra\textquotedblright{} + \emph{kâri} \textquotedblleft{}working\textquotedblright{}. \emph{Ezâfe kâri kardan} — to work overtime.

\begin{cousinweb}
\textbf{\fa{اضافه} \fa{کاری}} joins \emph{ezâfe}, from Arabic \emph{iḍāfa} \textquotedblleft{}addition, adjunct\textquotedblright{} (EtymDB), to \emph{kâri}, from the inherited Persian \emph{kâr} \textquotedblleft{}work\textquotedblright{} (EtymDB traces it to Middle Persian): added work.
\end{cousinweb}

\begin{grammarlens}[title={What you've built}]
One more word for telling, explaining and describing at B1.
\end{grammarlens}

\subsection*{Your turn}
\begin{itemize}
\raggedright
  \item \emph{Say it:} \emph{ezâfe kâri}
  \item \emph{Say it:} \emph{ezâfe kâri}, once more
  \item \emph{Say it:} a sentence of your own with \emph{ezâfe kâri}
  \item \emph{Recall:} say the Persian for a wage, then the Persian for an income, then say \emph{ezâfe kâri} again
\end{itemize}

\begin{tcolorbox}[breakable,colback=teal!4,colframe=teal!35!black,title={Before you move on}]
What does \fa{اضافه} \fa{کاری} mean? (\textquotedblleft{}Overtime\textquotedblright{}.) Say it once more.
\end{tcolorbox}
\section[\texorpdfstring{\fa{پاداش} (pâdâsh)}{پاداش (pâdâsh)}]{\fa{پاداش} (pâdâsh) — a bonus, a reward}

\label{lesson:FA-C378-padash}



\begin{quote}
Before the new one: say the Persian for an income, then the Persian for overtime.
\end{quote}

\subsection*{You'll want to know: \fa{پاداش}}
\textbf{\fa{پاداش}} (\emph{pâdâsh}) — \textquotedblleft{}a bonus, a reward\textquotedblright{}.

\begin{quote}
\fa{آخر} \fa{سال} \fa{همه} \fa{کارمندها} \fa{پاداش} \fa{گرفتند}.
\end{quote}

\begin{quote}
\emph{âkhar-e sâl hame-ye kârmandhâ pâdâsh gereftand.}
\end{quote}

\begin{quote}
At the end of the year all the employees got a bonus.
\end{quote}

Also a reward in a story; its opposite is \emph{mojâzât}, a punishment.

\begin{cousinweb}
\textbf{\fa{پاداش}} is inherited: EtymDB traces it to Middle Persian \emph{pādāšn} \textquotedblleft{}reward, retribution\textquotedblright{}, what is given back for a deed.
\end{cousinweb}

\begin{grammarlens}[title={What you've built}]
One more word for telling, explaining and describing at B1.
\end{grammarlens}

\subsection*{Your turn}
\begin{itemize}
\raggedright
  \item \emph{Say it:} \emph{pâdâsh}
  \item \emph{Say it:} \emph{pâdâsh}, once more
  \item \emph{Say it:} a sentence of your own with \emph{pâdâsh}
  \item \emph{Recall:} say the Persian for an income, then the Persian for overtime, then say \emph{pâdâsh} again
\end{itemize}

\begin{tcolorbox}[breakable,colback=teal!4,colframe=teal!35!black,title={Before you move on}]
What does \fa{پاداش} mean? (\textquotedblleft{}A bonus, a reward\textquotedblright{}.) Say it once more.
\end{tcolorbox}
\section[\texorpdfstring{\fa{مزایا} (mazâyâ)}{مزایا (mazâyâ)}]{\fa{مزایا} (mazâyâ) — benefits, perks}

\label{lesson:FA-C378-mazaya}



\begin{quote}
Before the new one: say the Persian for overtime, then the Persian for a bonus.
\end{quote}

\subsection*{You'll want to know: \fa{مزایا}}
\textbf{\fa{مزایا}} (\emph{mazâyâ}) — \textquotedblleft{}benefits, perks\textquotedblright{}.

\begin{quote}
\fa{این} \fa{شغل} \fa{حقوق} \fa{خوب} \fa{و} \fa{مزایای} \fa{زیادی} \fa{دارد}.
\end{quote}

\begin{quote}
\emph{in shoghl hoquq-e khub va mazâyâ-ye ziyâdi dârad.}
\end{quote}

\begin{quote}
This job has a good salary and a lot of benefits.
\end{quote}

The Arabic broken plural of \emph{maziyat} \textquotedblleft{}advantage\textquotedblright{}; at work, insurance, bonuses and leave.

\begin{cousinweb}
\textbf{\fa{مزایا}} is usually traced to Arabic \emph{mazāyā}, the broken plural of \emph{maziyya} \textquotedblleft{}merit, advantage\textquotedblright{}, the word behind \emph{maziyat}.
\end{cousinweb}

\begin{grammarlens}[title={What you've built}]
One more word for telling, explaining and describing at B1.
\end{grammarlens}

\subsection*{Your turn}
\begin{itemize}
\raggedright
  \item \emph{Say it:} \emph{mazâyâ}
  \item \emph{Say it:} \emph{mazâyâ}, once more
  \item \emph{Say it:} a sentence of your own with \emph{mazâyâ}
  \item \emph{Recall:} say the Persian for overtime, then the Persian for a bonus, then say \emph{mazâyâ} again
\end{itemize}

\begin{tcolorbox}[breakable,colback=teal!4,colframe=teal!35!black,title={Before you move on}]
What does \fa{مزایا} mean? (\textquotedblleft{}Benefits, perks\textquotedblright{}.) Say it once more.
\end{tcolorbox}
